\chapter{Chapter 5: Asymptotic Theory and Robust Inference}

\medskip\hrule\medskip

\section{Chapter Overview}

Chapter 4 derived exact finite-sample results under the strong assumption that errors are normally distributed. In practice, economic data are rarely normally distributed: wages are right-skewed, binary outcomes are Bernoulli, error variances vary across observations. This chapter develops \textbf{asymptotic (large-sample) theory} --- statistical results that hold approximately when $n$ is large, without requiring normality or homoskedasticity.

We derive the \textbf{asymptotic distribution of OLS}, show how to estimate the asymptotic variance under \textbf{heteroskedasticity} (non-constant variance) and \textbf{within-cluster correlation}, and provide practical guidance on when and how to use robust standard errors. These are the standard errors used in virtually all applied economic research today.

\medskip\hrule\medskip

\section{Learning Objectives}

After completing this chapter, students will be able to:

\begin{enumerate}
\item Explain why asymptotic theory is needed in applied econometrics.
\item State and apply the \textbf{Law of Large Numbers (LLN)} and \textbf{Central Limit Theorem (CLT)} to OLS.
\item Prove \textbf{consistency} of OLS under weaker conditions than Gauss-Markov.
\item Derive the \textbf{asymptotic distribution} of the OLS estimator.
\item Explain what \textbf{heteroskedasticity} is and why it invalidates classical standard errors.
\item Compute and interpret \textbf{heteroskedasticity-robust (HC) standard errors}.
\item Explain \textbf{cluster-robust standard errors} and when clustering is necessary.
\item Provide practical guidance on standard error choice in applied work.
\end{enumerate}

\medskip\hrule\medskip

\section{5.1 Why Asymptotic Theory?}

The Gauss-Markov assumptions, particularly normality (Assumption 4.5), are often implausible:

\begin{itemize}
\item \textbf{Discrete outcomes:} Earnings rounded to the nearest dollar, binary indicators of employment, counts of accidents. Error distributions are clearly non-normal.
\item \textbf{Non-spherical errors:} Wages vary more for high earners (heteroskedasticity); students in the same class share correlated outcomes (clustering); annual firm data have autocorrelated shocks.
\item \textbf{Misspecified functional form:} When the CEF is non-linear, OLS residuals systematically deviate from normality.
\end{itemize}

When the classical assumptions fail, we need \textbf{asymptotic theory} --- results that hold as $n \to \infty$ under weaker conditions. The key tools are the \textbf{Law of Large Numbers} (sample averages converge to population means) and the \textbf{Central Limit Theorem} (standardized sample averages are approximately normal for large $n$).

\medskip\hrule\medskip

\section{5.2 Consistency of OLS}

\textbf{Definition 5.1 (Consistency):} An estimator $\hat{\boldsymbol{\beta}}_n$ is \textbf{consistent} for $\boldsymbol{\beta}$ if $\hat{\boldsymbol{\beta}}_n \xrightarrow{p} \boldsymbol{\beta}$ as $n \to \infty$. That is, for any $\epsilon > 0$:
\[
\lim_{n \to \infty} P(\|\hat{\boldsymbol{\beta}}_n - \boldsymbol{\beta}\| > \epsilon) = 0
\]

\textbf{Theorem 5.1 (Consistency of OLS):} Suppose the following conditions hold:
\begin{enumerate}
\item The model is linear: $y_i = \mathbf{x}_i^T \boldsymbol{\beta} + u_i$.
\item The data $\{(y_i, \mathbf{x}_i)\}$ are i.i.d.
\item \textbf{Predeterminedness:} $E[\mathbf{x}_i u_i] = \mathbf{0}$ (orthogonality condition).
\item $E[\mathbf{x}_i \mathbf{x}_i^T]$ is finite and non-singular.
\end{enumerate}

Then $\hat{\boldsymbol{\beta}}^{OLS} \xrightarrow{p} \boldsymbol{\beta}$.

\textbf{Proof:}

\[
\hat{\boldsymbol{\beta}} = \boldsymbol{\beta} + \left(\frac{1}{n}\mathbf{X}^T\mathbf{X}\right)^{-1} \left(\frac{1}{n}\mathbf{X}^T\mathbf{u}\right)
\]

By the LLN:
\[
\frac{1}{n}\mathbf{X}^T\mathbf{X} = \frac{1}{n}\sum_{i=1}^n \mathbf{x}_i \mathbf{x}_i^T \xrightarrow{p} E[\mathbf{x}_i \mathbf{x}_i^T] \equiv \mathbf{S}_{XX}
\]

\[
\frac{1}{n}\mathbf{X}^T\mathbf{u} = \frac{1}{n}\sum_{i=1}^n \mathbf{x}_i u_i \xrightarrow{p} E[\mathbf{x}_i u_i] = \mathbf{0}
\]

By Slutsky's theorem and the continuous mapping theorem:

\[
\hat{\boldsymbol{\beta}} \xrightarrow{p} \boldsymbol{\beta} + \mathbf{S}_{XX}^{-1} \mathbf{0} = \boldsymbol{\beta} \quad \square
\]

\subsection{5.2.1 Consistency vs. Unbiasedness}

\textbf{Key distinction:} Consistency only requires the weaker condition $E[\mathbf{x}_i u_i] = \mathbf{0}$ (predeterminedness / zero covariance), not strict exogeneity $E[u_i \mid \mathbf{X}] = \mathbf{0}$.

\begin{itemize}
\item Strict exogeneity $\Rightarrow$ predeterminedness, but NOT vice versa.
\item Predeterminedness is enough for consistency; strict exogeneity is needed for unbiasedness.
\end{itemize}

This means OLS can be \textbf{consistent but biased in small samples} --- the bias vanishes as $n \to \infty$.

\textbf{Practical implication:} In many panel data models (with lagged dependent variables), OLS is biased in small samples but consistent in large samples. For small panels, this bias can be substantial.

\medskip\hrule\medskip

\section{5.3 Asymptotic Normality and the Sandwich Variance}

\textbf{Theorem 5.2 (Asymptotic Normality of OLS):} Under the i.i.d. assumptions of Theorem 5.1, plus $E[\mathbf{x}_i \mathbf{x}_i^T u_i^2]$ finite:

\[
\sqrt{n}(\hat{\boldsymbol{\beta}} - \boldsymbol{\beta}) \xrightarrow{d} N\!\left(\mathbf{0}, \mathbf{S}_{XX}^{-1} \boldsymbol{\Omega} \mathbf{S}_{XX}^{-1}\right)
\]

where $\mathbf{S}_{XX} = E[\mathbf{x}_i \mathbf{x}_i^T]$ and $\boldsymbol{\Omega} = E[\mathbf{x}_i \mathbf{x}_i^T u_i^2] = E[u_i^2 \mathbf{x}_i \mathbf{x}_i^T]$.

\textbf{Proof:} The key step is applying the \textbf{Multivariate Central Limit Theorem (MCLT)} to $\frac{1}{\sqrt{n}}\mathbf{X}^T\mathbf{u}$:

\[
\frac{1}{\sqrt{n}}\mathbf{X}^T\mathbf{u} = \frac{1}{\sqrt{n}}\sum_{i=1}^n \mathbf{x}_i u_i \xrightarrow{d} N(\mathbf{0}, \boldsymbol{\Omega})
\]

By the LLN, $\frac{1}{n}\mathbf{X}^T\mathbf{X} \xrightarrow{p} \mathbf{S}_{XX}$. By Slutsky's theorem:

\[
\sqrt{n}(\hat{\boldsymbol{\beta}} - \boldsymbol{\beta}) = \left(\frac{1}{n}\mathbf{X}^T\mathbf{X}\right)^{-1} \frac{1}{\sqrt{n}}\mathbf{X}^T\mathbf{u} \xrightarrow{d} \mathbf{S}_{XX}^{-1} \cdot N(\mathbf{0}, \boldsymbol{\Omega}) = N(\mathbf{0}, \mathbf{S}_{XX}^{-1}\boldsymbol{\Omega}\mathbf{S}_{XX}^{-1}) \quad \square
\]

\begin{figure}[htbp]
  \centering
  \begin{tikzpicture}
    \begin{axis}[
      width=11cm, height=6.5cm,
      xlabel={$\sqrt{n}(\hat{\beta} - \beta_0)$},
      ylabel={Density},
      xmin=-4, xmax=4,
      ymin=0, ymax=0.5,
      axis lines=left,
      tick style={draw=none},
      xtick={0},
      xticklabels={0},
      ytick=\empty,
      legend style={at={(0.97,0.95)}, anchor=north east, font=\small},
    ]
      % Small n distribution (wider, less normal)
      \addplot[thick, gray!60, domain=-4:4, samples=100]
        {0.38*exp(-0.5*(x)^2/1.8^2)/sqrt(2*pi*1.8^2)*4};

      % Medium n
      \addplot[thick, econblue!60, domain=-4:4, samples=100]
        {0.4*exp(-0.5*(x)^2/1.2^2)/sqrt(2*pi*1.2^2)*4};

      % Large n (approaching normal)
      \addplot[very thick, econblue, domain=-4:4, samples=100]
        {0.42*exp(-0.5*(x)^2)/sqrt(2*pi)*4};

      % Standard normal (limiting)
      \addplot[thick, red!70!black, dashed, domain=-4:4, samples=100]
        {exp(-0.5*x^2)/sqrt(2*pi)};

      \addlegendentry{Small $n$}
      \addlegendentry{Moderate $n$}
      \addlegendentry{Large $n$}
      \addlegendentry{$\mathcal{N}(0,V_\beta)$ (limiting)}

      \node[font=\footnotesize, econblue] at (axis cs:-2.5, 0.38) {$\sqrt{n}(\hat{\beta}-\beta_0)\xrightarrow{d}\mathcal{N}(0,V_\beta)$};
    \end{axis}
  \end{tikzpicture}
  \caption{Asymptotic normality of the OLS estimator. As the sample size $n$ increases, the scaled and centered OLS estimator $\sqrt{n}(\hat{\beta} - \beta_0)$ converges in distribution to a normal random variable. The limiting distribution $\mathcal{N}(0, V_\beta)$ (dashed red) provides the basis for inference in large samples.}
  \label{fig:ch05_asymptotic_normality}
\end{figure}

\subsection{5.3.1 The Sandwich Variance Estimator}

The matrix $\mathbf{S}_{XX}^{-1}\boldsymbol{\Omega}\mathbf{S}_{XX}^{-1}$ is called the \textbf{sandwich variance} (or \textbf{Huber-White} variance). In finite samples:

\[
\text{Var}(\hat{\boldsymbol{\beta}}) \approx \frac{1}{n} \mathbf{S}_{XX}^{-1} \boldsymbol{\Omega} \mathbf{S}_{XX}^{-1} = (\mathbf{X}^T\mathbf{X})^{-1}\left[\mathbf{X}^T\boldsymbol{\Omega}\mathbf{X}\right](\mathbf{X}^T\mathbf{X})^{-1}
\]

\textbf{Under homoskedasticity:} $E[u_i^2 \mid \mathbf{x}_i] = \sigma^2$, so $\boldsymbol{\Omega} = \sigma^2 \mathbf{S}_{XX}$, and the sandwich variance simplifies to:

\[
\frac{\sigma^2}{n} \mathbf{S}_{XX}^{-1} \approx \sigma^2 (\mathbf{X}^T\mathbf{X})^{-1}
\]

This is the classical variance formula --- consistent with the Gauss-Markov result.

\textbf{Under heteroskedasticity:} $E[u_i^2 \mid \mathbf{x}_i] = \sigma_i^2 \neq \sigma^2$, the classical formula gives the wrong variance. We need to estimate $\boldsymbol{\Omega}$ consistently.

\medskip\hrule\medskip

\section{5.4 Heteroskedasticity-Robust Standard Errors}

\textbf{Definition 5.2 (Heteroskedasticity):} Errors are heteroskedastic if $V(u_i \mid \mathbf{x}_i)$ varies across observations.

\textbf{Example 5.1:} In a wage regression, it is common for high earners to have greater dispersion in wages than low earners --- the variance of $u_i$ may be an increasing function of predicted wages.

\subsection{5.4.1 Consequences of Ignoring Heteroskedasticity}

If we use the classical variance formula $\hat{\sigma}^2(\mathbf{X}^T\mathbf{X})^{-1}$ under heteroskedasticity:
\begin{itemize}
\item OLS is still \textbf{consistent} (consistency does not require spherical errors).
\item But standard errors are \textbf{wrong}: they may be too large or too small.
\item $t$-tests and $F$-tests are \textbf{invalid}: rejection rates under the null are not $\alpha$.
\end{itemize}

\subsection{5.4.2 The Heteroskedasticity-Robust (HC) Estimator}

\textbf{White (1980)} proposed a consistent estimator of the sandwich variance that does not require spherical errors:

\[
\widehat{\text{Var}}_{HC}(\hat{\boldsymbol{\beta}}) = (\mathbf{X}^T\mathbf{X})^{-1} \left(\sum_{i=1}^n \hat{u}_i^2 \mathbf{x}_i \mathbf{x}_i^T\right) (\mathbf{X}^T\mathbf{X})^{-1}
\]

This is a sample analogue of the sandwich formula, replacing $E[u_i^2 \mathbf{x}_i \mathbf{x}_i^T]$ with $\frac{1}{n}\sum_{i=1}^n \hat{u}_i^2 \mathbf{x}_i \mathbf{x}_i^T$.

\textbf{Why does this work?} By the consistency of OLS, $\hat{u}_i \approx u_i$ for large $n$. So $\hat{u}_i^2 \mathbf{x}_i \mathbf{x}_i^T$ is a consistent estimator of $E[u_i^2 \mathbf{x}_i \mathbf{x}_i^T]$.

\subsection{5.4.3 Small-Sample Corrections (HC1, HC2, HC3)}

In small samples, $\hat{u}_i^2$ is a biased downward estimate of $u_i^2$ (because OLS fits the residuals). Several corrections are used in practice:

\begin{center}
\begin{tabular}{@{} l l l @{}}
\toprule
\textbf{Version} & \textbf{Formula} & \textbf{Motivation} \\
\midrule
HC0 & $\hat{u}_i^2$ & Original White (1980) estimator \\
HC1 & $\frac{n}{n-k}\hat{u}_i^2$ & Degrees-of-freedom correction \\
HC2 & $\frac{\hat{u}_i^2}{1 - h_{ii}}$ & Corrects for leverage \\
HC3 & $\frac{\hat{u}_i^2}{(1 - h_{ii})^2}$ & Better in very small samples \\
\bottomrule
\end{tabular}
\end{center}

where $h_{ii} = [\mathbf{P_X}]_{ii}$ is the $i$-th diagonal of the hat matrix (leverage of observation $i$).

\textbf{In practice:} Use HC3 for small samples ($n < 250$), HC1 for larger samples. The \texttt{sandwich} and \texttt{lmtest} packages in R provide all variants.

\begin{keybox}{The Rule of Thumb:** In modern applied econometrics, always report heteroskedasticity-robust standard errors unless you have a strong theoretical reason to believe errors are homoskedastic.}
\end{keybox}

\medskip\hrule\medskip

\section{5.5 Cluster-Robust Standard Errors}

\subsection{5.5.1 When Clustering Matters}

\textbf{Clustering arises when observations are not independent.} Consider:

\[
y_{ig} = \mathbf{x}_{ig}^T \boldsymbol{\beta} + u_{ig}
\]

where $i$ indexes individuals within group (cluster) $g$. If individuals within the same group share common unobservables (e.g., students in the same school share a teacher quality shock), then $\text{Cov}(u_{ig}, u_{jg}) \neq 0$ for $i \neq j$.

\textbf{The variance-covariance structure is block-diagonal:}

\[
E[\mathbf{u}\mathbf{u}^T] = \begin{pmatrix} \boldsymbol{\Sigma}_1 & \mathbf{0} & \cdots & \mathbf{0} \\ \mathbf{0} & \boldsymbol{\Sigma}_2 & \cdots & \mathbf{0} \\ \vdots & & \ddots & \vdots \\ \mathbf{0} & \mathbf{0} & \cdots & \boldsymbol{\Sigma}_G \end{pmatrix}
\]

where $\boldsymbol{\Sigma}_g = E[\mathbf{u}_g \mathbf{u}_g^T]$ allows arbitrary within-cluster correlation.

\subsection{5.5.2 The Cluster-Robust Sandwich Estimator}

\textbf{Definition 5.3 (Cluster-Robust SE):}

\[
\widehat{\text{Var}}_{CR}(\hat{\boldsymbol{\beta}}) = (\mathbf{X}^T\mathbf{X})^{-1} \left(\sum_{g=1}^G \mathbf{X}_g^T \hat{\mathbf{u}}_g \hat{\mathbf{u}}_g^T \mathbf{X}_g\right) (\mathbf{X}^T\mathbf{X})^{-1}
\]

where $\mathbf{X}_g$ and $\hat{\mathbf{u}}_g$ are the design matrix and residuals for cluster $g$.

\textbf{Key difference from HC:} HC sums over individual observations; CR sums over clusters. Within each cluster, we allow arbitrary covariance in residuals.

\subsection{5.5.3 Practical Guidance on When to Cluster}

Clustering is needed when:

\begin{enumerate}
\item \textbf{Treatment is at a higher level than outcomes.} E.g., a minimum wage law applies to a state, but outcomes are individual workers. Cluster at the state level.
\end{enumerate}

\begin{enumerate}
\item \textbf{Panel data with repeated observations.} Each individual appears multiple times; errors are autocorrelated over time. Cluster at the individual level.
\end{enumerate}

\begin{enumerate}
\item \textbf{Assignment of treatment is clustered.} E.g., a program assigns villages to treatment; outcomes are individual household members. Cluster at the village level.
\end{enumerate}

\textbf{Rule of thumb:} Cluster at the level at which treatment varies (or at which observations are correlated).

\textbf{Warning --- few clusters:} Cluster-robust SEs rely on $G \to \infty$. With few clusters ($G < 20$), they are unreliable. In such cases, use the \textbf{wild cluster bootstrap} (see Cameron and Miller, 2015).

\textbf{Example 5.2 (Card and Krueger, 1994):} In their study of the effect of a New Jersey minimum wage increase on fast-food employment, Card and Krueger cluster at the restaurant level --- each restaurant is surveyed before and after the policy change, so two observations per restaurant are correlated.

\medskip\hrule\medskip

\section{5.6 When to Cluster? Practical Guidance}

A common question from students: "Should I cluster?" The answer depends on the data structure, not the size of the coefficient:

\begin{center}
\begin{tabular}{@{} l l l @{}}
\toprule
\textbf{Situation} & \textbf{Cluster?} & \textbf{Level} \\
\midrule
Random sample, i.i.d. observations & No & --- \\
Survey with multi-stage sampling & Yes & Primary sampling unit \\
Difference-in-differences & Yes & State/region/treatment unit \\
Panel data & Yes & Individual/firm \\
Cross-section with geographic variation & Often & Geographic unit \\
RCT with cluster randomization & Yes & Randomization unit \\
\bottomrule
\end{tabular}
\end{center}

\textbf{Never} choose whether to cluster based on whether it changes your results. Cluster based on the structure of the data.

\medskip\hrule\medskip

\section{5.7 Empirical Illustration: Wages and Education Revisited}

Consider the CPS wage regression:

\[
\log(\text{wage}_i) = \beta_1 + \beta_2 \cdot \text{educ}_i + \beta_3 \cdot \text{exp}_i + \beta_4 \cdot \text{exp}_i^2 + u_i
\]

\textbf{Table 5.1: OLS Estimates with Different Standard Errors}

\begin{center}
\begin{tabular}{@{} l l l l @{}}
\toprule
\textbf{} & \textbf{Classical SE} & \textbf{HC1 SE} & \textbf{HC3 SE} \\
\midrule
Education & 0.0709 (0.0010) & 0.0709 (0.0011) & 0.0709 (0.0011) \\
Experience & 0.0343 (0.0010) & 0.0343 (0.0013) & 0.0343 (0.0013) \\
Exp$^{2}$ & $-$0.0005 (0.000) & $-$0.0005 (0.000) & $-$0.0005 (0.000) \\
\bottomrule
\end{tabular}
\end{center}

Standard errors in parentheses.

In this large sample, classical and robust SEs are similar because $n$ is large. The key point: in datasets with clear heteroskedasticity patterns, or with clustering, the differences can be substantial.

\medskip\hrule\medskip

\section{Chapter Summary}

\begin{itemize}
\item Asymptotic theory delivers results about $\hat{\boldsymbol{\beta}}$ under weaker conditions than finite-sample theory.
\item OLS is \textbf{consistent} under predeterminedness ($E[\mathbf{x}_i u_i] = \mathbf{0}$) --- weaker than strict exogeneity.
\item The \textbf{asymptotic distribution} of $\sqrt{n}(\hat{\boldsymbol{\beta}} - \boldsymbol{\beta})$ is normal with sandwich variance $\mathbf{S}_{XX}^{-1}\boldsymbol{\Omega}\mathbf{S}_{XX}^{-1}$.
\item \textbf{Heteroskedasticity} means $V(u_i \mid \mathbf{x}_i)$ is not constant. It does not bias OLS but invalidates classical SEs.
\item \textbf{White (HC) standard errors} provide consistent inference under heteroskedasticity of unknown form.
\item \textbf{Cluster-robust standard errors} are needed when observations within groups are correlated; they cluster the "meat" of the sandwich at the group level.
\item In modern applied work, \textbf{always report robust standard errors} unless there is a strong reason to believe errors are spherical.
\end{itemize}

\medskip\hrule\medskip

\section{Key Terms}

\textbf{Consistency} --- $\hat{\boldsymbol{\beta}} \xrightarrow{p} \boldsymbol{\beta}$ as $n \to \infty$.

\textbf{Predeterminedness} --- $E[\mathbf{x}_i u_i] = \mathbf{0}$; weaker than strict exogeneity.

\textbf{Asymptotic normality} --- $\sqrt{n}(\hat{\boldsymbol{\beta}} - \boldsymbol{\beta}) \xrightarrow{d} N(\mathbf{0}, \mathbf{V})$ for some variance matrix $\mathbf{V}$.

\textbf{Sandwich variance} --- $\mathbf{S}_{XX}^{-1}\boldsymbol{\Omega}\mathbf{S}_{XX}^{-1}$; the asymptotic variance of OLS.

\textbf{Heteroskedasticity} --- non-constant error variance: $V(u_i \mid \mathbf{x}_i) \neq \sigma^2$.

\textbf{Heteroskedasticity-robust (HC) SE} --- consistent standard errors that allow for heteroskedasticity.

\textbf{Cluster-robust SE} --- consistent standard errors that allow for within-cluster correlation.

\textbf{Wild cluster bootstrap} --- a resampling method for inference with few clusters.

\medskip\hrule\medskip

\section{Exercises}

\subsection{Conceptual Questions}

\textbf{5.1} Explain the difference between \textbf{unbiasedness} and \textbf{consistency}. Give an example of an estimator that is consistent but biased in small samples.

\textbf{5.2} Why does heteroskedasticity not cause OLS to be biased or inconsistent? What does it affect?

\textbf{5.3} A researcher runs a difference-in-differences regression using state-level panel data (50 states, 10 years). A colleague says "You have 500 observations, so clustering doesn't matter." Evaluate this claim.

\textbf{5.4} What is the "sandwich" structure of the heteroskedasticity-robust variance estimator? Why is it called a sandwich?

\textbf{5.5} You are analyzing school data where 30 students are observed in each of 50 schools. Should you cluster by student, school, or not at all? Justify your answer.

\subsection{Analytical Questions}

\textbf{5.6} Prove that under predeterminedness and i.i.d. sampling, OLS is consistent.

\textbf{5.7} Suppose $V(u_i \mid x_i) = \sigma^2 x_i^2$. Derive the true variance of OLS in the bivariate regression of $y_i$ on $x_i$. Show it differs from the classical formula $\sigma^2/\text{SST}_x$.

\textbf{5.8} Show that when $V(u_i \mid \mathbf{x}_i) = \sigma^2$ (homoskedasticity), the HC0 sandwich estimator is NOT equal to the classical estimator, but converges to the same limit as $n \to \infty$.

\textbf{5.9} Consider panel data $y_{it} = \mathbf{x}_{it}^T\boldsymbol{\beta} + \alpha_i + \epsilon_{it}$ where $\alpha_i$ is an individual-specific shock and $\epsilon_{it}$ is i.i.d. If you run OLS pooling all observations without fixed effects, $\hat{\boldsymbol{\beta}}$ is generally biased. But even if you control for fixed effects (making $\hat{\boldsymbol{\beta}}$ consistent), why might you still need to cluster standard errors at the individual level?

\subsection{Applied Questions}

\textbf{5.10} Using the CPS data, estimate the wage regression with (a) classical SEs, (b) HC1 SEs, and (c) HC3 SEs. Report all three and compare. When do they differ?

\textbf{5.11} In R, demonstrate via Monte Carlo simulation that the classical $t$-test has the wrong size (rejection rate $\neq 5\%$ under $H_0$) when errors are heteroskedastic, but the HC-robust $t$-test maintains the correct size.

\medskip\hrule\medskip

\section{R Lab 5: Robust Standard Errors in Practice}

\subsection{Setup}

\begin{lstlisting}
library(tidyverse)
library(sandwich)   # for vcovHC
library(lmtest)    # for coeftest
library(AER)

data("CPS1988")
CPS1988 <- CPS1988 %>% mutate(log_wage = log(wage))

model <- lm(log_wage ~ education + experience + I(experience^2), data = CPS1988)
\end{lstlisting}

\subsection{Exercise 5.1: Compare Standard Errors}

\begin{lstlisting}
# Classical SEs
coeftest(model)

# HC0 (White 1980)
coeftest(model, vcov = vcovHC(model, type = "HC0"))

# HC1 (degrees-of-freedom correction)
coeftest(model, vcov = vcovHC(model, type = "HC1"))

# HC3 (best for small samples)
coeftest(model, vcov = vcovHC(model, type = "HC3"))
\end{lstlisting}

\subsection{Exercise 5.2: Detecting Heteroskedasticity}

\begin{lstlisting}
# Residuals vs fitted values plot
df_resid <- data.frame(
  fitted   = fitted(model),
  residuals = residuals(model)
)

ggplot(df_resid, aes(x = fitted, y = residuals)) +
  geom_point(alpha = 0.2, color = "steelblue") +
  geom_hline(yintercept = 0, color = "red") +
  geom_smooth(method = "loess", se = FALSE, color = "darkred") +
  labs(
    title = "Residuals vs. Fitted Values",
    subtitle = "Pattern indicates heteroskedasticity",
    x = "Fitted Values", y = "Residuals"
  ) +
  theme_minimal()

# Breusch-Pagan test for heteroskedasticity
library(lmtest)
bptest(model)
\end{lstlisting}

\subsection{Exercise 5.3: Cluster-Robust Standard Errors}

\begin{lstlisting}
library(sandwich)

# Simulate panel-like clustering by state (using ethnicity as a proxy cluster)
# In real applications, cluster = geographic unit or treatment unit

# Cluster by 'ethnicity' (as illustrative proxy)
vcov_cluster <- vcovCL(model, cluster = ~ CPS1988$ethnicity)
coeftest(model, vcov = vcov_cluster)
\end{lstlisting}

\subsection{Exercise 5.4: Monte Carlo --- Size of Classical vs. Robust t-tests under Heteroskedasticity}

\begin{lstlisting}
set.seed(42)
n_sims <- 5000
n      <- 200
reject_classical <- numeric(n_sims)
reject_robust    <- numeric(n_sims)

for (s in 1:n_sims) {
  x <- rnorm(n)
  u <- rnorm(n, sd = abs(x))   # heteroskedastic: V(u|x) = x^2
  y <- 1 + 0 * x + u            # H0: beta = 0 is TRUE

  fit <- lm(y ~ x)
  
  # Classical t-stat
  t_classic <- summary(fit)$coefficients["x", "t value"]
  reject_classical[s] <- abs(t_classic) > 1.96
  
  # Robust t-stat
  t_robust <- coeftest(fit, vcov = vcovHC(fit, "HC1"))["x", "t value"]
  reject_robust[s] <- abs(t_robust) > 1.96
}

cat("Classical test rejection rate:", round(mean(reject_classical), 3),
    "(should be ~0.05)\n")
cat("Robust test rejection rate:   ", round(mean(reject_robust), 3),
    "(should be ~0.05)\n")
\end{lstlisting}

\textbf{Expected finding:} The classical test rejects too often (size > 5\%) under heteroskedasticity; the robust test maintains approximately the correct size.

\medskip\hrule\medskip

\textit{End of Chapter 5}

\medskip\hrule\medskip

\textbf{References}

White, H. (1980). A heteroskedasticity-consistent covariance matrix estimator and a direct test for heteroskedasticity. \textit{Econometrica}, 48(4), 817--838.

Davidson, R., and MacKinnon, J. G. (2004). \textit{Econometric Theory and Methods}. Chapter 3.

Angrist, J. D., and Pischke, J. (2009). \textit{Mostly Harmless Econometrics}. Chapter 8.

Cameron, A. C., and Miller, D. L. (2015). A practitioner's guide to cluster-robust inference. \textit{Journal of Human Resources}, 50(2), 317--372.

Cunningham, S. (2021). \textit{Causal Inference: The Mixtape}. Chapter on regression.
